\section{System-by-System Analysis}
\label{sec:systems}

We organize systems into seven categories: production systems
(\S\ref{sec:production}), standalone memory servers
(\S\ref{sec:standalone}), forgetting-aware systems
(\S\ref{sec:forgetting-systems}), framework-level memory
(\S\ref{sec:framework}), coding agent memory (\S\ref{sec:coding-agents}),
file-based memory (\S\ref{sec:filebased}), and research architectures
(\S\ref{sec:research}). For each system we identify the memory
types supported (mapped to the CoALA taxonomy~\cite{sumers2024coala}),
the retrieval architecture (mapped to the RAG
taxonomy~\cite{gao2023rag}), the lifecycle operations, and the
infrastructure requirements. Quantitative cross-comparisons appear in
\S\ref{sec:comparison}; here we focus on what makes each system
architecturally distinct.

%% ===============================================================
\subsection{Production Systems}
\label{sec:production}

\paragraph{Mem0.}
Mem0~\cite{choudhary2025mem0} implements a tiered memory architecture
backed by vector databases (Qdrant, Pinecone, or Chroma), a relational
store for metadata, and a Neo4j graph layer.  With v1.0.0 (early 2026),
the graph component---branded \emph{Mem0-g}---became a first-class
subsystem: an entity extraction pipeline identifies entities and
relationships from incoming interactions, a relation generation stage
links them into a knowledge graph, and a contradiction resolution module
detects and reconciles conflicting facts.  Memory formation remains
LLM-dependent: every incoming interaction is routed through an LLM that
extracts facts, user preferences, and personality traits, then decides
whether to create, update, or deduplicate entries.  Retrieval combines
vector similarity with graph traversal and contextual priority scoring,
placing Mem0 in the Advanced RAG category.

Mem0 v1.0.0 introduced three lifecycle mechanisms absent from the
original release: (1)~\emph{dynamic forgetting/decay} that prunes
low-relevance entries based on adaptive decay curves rather than fixed
TTLs, (2)~\emph{TTL with expiration dates} enabling time-bounded
memories that auto-expire, and (3)~a \emph{Memory Compression Engine}
that merges and consolidates redundant facts into compact
representations.  These additions make Mem0 one of the few production
systems with both forgetting and consolidation.  The authors report a
26\% accuracy improvement over OpenAI's built-in memory on the LoCoMo
benchmark and 91\% lower p95 latency versus full-context approaches.
With 41K GitHub stars, 14M+ downloads, and selection as the exclusive
memory provider for the AWS Agent SDK, Mem0 is the most widely adopted
standalone memory library.

The key architectural tradeoff remains Mem0's heavy infrastructure
requirement (2--3 external services) and its reliance on an LLM for
every memory write.  This makes Mem0 well-suited for cloud-deployed,
multi-user SaaS products with dedicated operations teams, but less
appropriate for privacy-first or offline use cases.  The full graph
features (Mem0-g with entity extraction, relation generation, and
contradiction resolution) require the \$249/month Pro tier, limiting
accessibility for individual developers and small teams.

\paragraph{Zep / Graphiti.}
Graphiti~\cite{dempsey2025graphiti}, the engine underlying Zep's memory
service, constructs a hierarchical temporal knowledge graph on Neo4j
with three subgraph layers: community (cluster summaries), entity
(semantic entities and relationships), and episodic (raw conversation
episodes). Retrieval uses a triple-hybrid pipeline---cosine similarity,
Okapi BM25, and breadth-first graph traversal---fused via Reciprocal
Rank Fusion (RRF) with subsequent maximal marginal relevance (MMR)
diversification. Multi-stage reranking considers episode mentions, node
distance, and cross-signal fusion. This places Zep/Graphiti firmly in
the Modular RAG category, alongside a small minority of systems
(17.9\%, Table~\ref{tab:rag-distribution}).

Zep's most distinctive lifecycle feature is \emph{bi-temporal
invalidation}: facts are never deleted, only marked with validity
periods recording when each fact was true in the real world and when
the system learned it. This conservative approach preserves full
provenance but means the knowledge graph grows monotonically. Entity
extraction requires an LLM (typically \texttt{gpt-4o-mini}), though
retrieval itself is LLM-free. The authors report 94.8\% accuracy on
the DMR benchmark and up to 18.5\% accuracy improvement on
LongMemEval.

\paragraph{Letta (formerly MemGPT).}
Letta~\cite{packer2023memgpt} introduced the operating-system metaphor
for agent memory: core memory serves as ``RAM'' (always in-context,
editable by the agent), recall memory stores conversation history, and
archival memory provides unbounded ``disk'' storage via a vector or
graph database. The defining architectural choice is that the LLM
\emph{is} the memory manager---it decides what to store, evict, and
retrieve via explicit tool calls (\texttt{core\_memory\_append},
\texttt{archival\_memory\_search}, etc.).

This agent-driven design makes Letta's memory quality entirely
dependent on LLM capability. When the context window fills, the
system evicts approximately 70\% of messages through summarization.
A 2025 extension adds ``sleep-time agents'' that handle asynchronous
reorganization of archival memory outside the main interaction loop.
Letta's CoALA mapping spans all three memory types (working, episodic,
semantic), but its retrieval is Naive RAG---agent-driven tool calls
with single-signal vector similarity. The OS metaphor proved
influential: MemoryOS~\cite{kang2025memoryos} and OpenViking both
adopted tiered hierarchies, though with different management
strategies.

\paragraph{Hindsight.}
Hindsight~\cite{latimer2025hindsight} introduces the most novel memory
decomposition in the production landscape: four logical networks for
\emph{world} (objective facts), \emph{experience} (agent
interactions), \emph{opinion} (subjective beliefs with confidence
scores), and \emph{observation} (preference-neutral entity summaries).
The opinion network is architecturally unique---each belief carries an
explicit confidence score that evolves as new evidence arrives,
enabling the system to distinguish between well-established facts and
tentative assessments.

Retrieval uses two specialized components: Tempr (temporal entity
memory priming) for time-aware recall and Cara (coherent adaptive
reasoning) for multi-step inference across networks. Hindsight does
not implement explicit forgetting; instead, belief confidence scores
decay or strengthen as evidence accumulates, producing an emergent
form of selective retention. The authors claim state-of-the-art
results: 91.4\% on LongMemEval and 89.61\% on LoCoMo. The four-network
decomposition is finer-grained than most systems' two- or three-store
architectures, but Hindsight lacks a graph overlay and does not
implement principled forgetting beyond confidence evolution.

\paragraph{ChatGPT Memory (OpenAI).}
OpenAI's built-in memory for ChatGPT~\cite{openai2024memoryfaq}
represents the largest-scale deployment of persistent agent memory,
available to hundreds of millions of users since early 2024. The
architecture implements a two-tier system: \emph{saved memories}
(explicit facts the user asks ChatGPT to remember, persisted
indefinitely until deleted) and \emph{reference chat history} (implicit
recall from past conversations, where details change over time as the
model decides what is most helpful to retain). Saved memories function
as auto-managed custom instructions---the model can create, update,
combine, or remove them without user intervention. When both tiers are
active, relevant information from both is injected as context into new
conversations.

Notably, ChatGPT Plus and Pro users have access to automatic memory
management that prioritizes memories based on recency and frequency of
mention, deprioritizing less-relevant memories to the
``background''---effectively a rudimentary forgetting mechanism based on
single-curve decay. The system also enforces a capacity limit (a
``memory full'' state), creating implicit pressure toward concise,
high-value memories. However, there is no structured retrieval:
memories are injected as context without ranking, no graph structure
links related memories, and no consolidation promotes episodic
interactions into semantic knowledge. Mem0 reports a 26\% accuracy
improvement over OpenAI Memory on LoCoMo~\cite{choudhary2025mem0},
consistent with our finding that full-context injection degrades as
memory stores grow. In our taxonomy, ChatGPT Memory is a Naive RAG
system with implicit preference storage and rudimentary forgetting, but
no graph, no consolidation, and no hybrid retrieval.

\paragraph{Claude Memory (Anthropic).}
Anthropic provides two distinct memory mechanisms. The \emph{Claude API
Memory Tool}~\cite{anthropic2026memorytool} enables developers to give
Claude persistent memory across conversations through a file-based
storage directory. Claude autonomously decides when to save, update, or
delete memories, storing them as structured Markdown documents organized
by topic. The tool exposes three operations---create, update, and
delete---and memories are injected as context at the start of each
conversation. This is architecturally simpler than Mem0 or Zep: no
vector retrieval, no graph traversal, no scoring---every saved memory
is included in full.

At the product level, \emph{Claude Code}'s auto-memory
system~\cite{gurgone2026claudememory} maintains a
\texttt{MEMORY.md} file per project
(\texttt{\~{}/.claude/\allowbreak{}projects/\allowbreak{}<path>/\allowbreak{}memory/\allowbreak{}MEMORY.md}), injected into
the system prompt at session start. A hard 200-line limit constrains
the file; beyond this, only the first 200 lines are loaded with a
warning nudging users toward a hub-and-spoke pattern where
\texttt{MEMORY.md} serves as an index and details live in separate
topic files (this topic-file loading is behind a feature flag and not
yet active by default). Custom agents have a three-scope memory
system: project-level (git-committable, team-shared), local
(machine-only), and user-level (global across projects). Both the API
tool and Claude Code's system are canonical examples of the
MEMORY.md problem analyzed in \S\ref{sec:memorymd}: full-context
injection with no retrieval, no lifecycle management, and a linear
scaling cost in memory size.

\paragraph{Gemini Memory (Google).}
Google's Gemini~\cite{gemini2026memory} deploys persistent memory across
its consumer assistant and Workspace applications (Gmail, Docs, Sheets,
Slides, Drive).  The architecture uses vector embeddings to store
``important user facts'' extracted from conversations, injected as
context in subsequent sessions.  A ``Recent chats'' panel (rolled out
March 2026) enables users to revisit and continue prior conversations
with full context preservation across Workspace apps.  Notably, Gemini
adopts an \emph{opt-out} privacy model---memory is enabled by default,
the opposite of ChatGPT's opt-in approach.  Memory management is
currently all-or-nothing: users can disable memory entirely or clear
all memories, but selective deletion of individual memories was not
yet available during testing.  Despite Gemini's industry-leading 1M-token
context window, the memory system remains Naive RAG: stored facts are
injected without ranking, graph structure, or lifecycle operations.

\paragraph{Grok Memory (xAI).}
xAI introduced memory for Grok in April 2025~\cite{grok2025memory},
enabling the chatbot to retain user preferences, names, and topical
interests across conversations on X (formerly Twitter), Grok.com, and
mobile apps.  The implementation follows the ChatGPT pattern: the model
decides what to remember, stored memories are transparent and
individually deletable, and a ``Personalize with Memories'' toggle
provides user control.  Architecturally, Grok Memory is Naive RAG with
implicit preference storage---no retrieval ranking, no graph, no
forgetting mechanism.  Its primary distinguishing feature is integration
with the X social platform, providing conversational context that other
assistants lack.

\paragraph{Meta AI Memory.}
Meta deploys memory across WhatsApp, Instagram, Messenger, and
meta.ai~\cite{meta2025aimemory}, potentially reaching the largest user
base of any platform memory system.  Memory stores names, preferences,
and recent topics across 1:1 chats, injected at runtime into the
model's context window.  Crucially, memory content \emph{consumes}
active context tokens---unlike systems that maintain a separate memory
store, Meta AI's injected memories reduce the available window for
conversation, creating implicit pressure toward concise memory as
histories grow.  The system provides no retrieval: all stored memories
are injected regardless of relevance, placing it firmly in the Naive
RAG category.

\paragraph{Perplexity Memory.}
Perplexity~\cite{perplexity2025memory} stands apart from the other
platform memory systems with two distinctive features.  First,
\emph{cross-model portability}: memory persists across every model
Perplexity offers (Claude, GPT, Gemini, and others), allowing users
to switch models between questions while retaining their full context
history---unique among all platforms surveyed.  Second,
\emph{citation-based transparency}: when memory influences a response,
Perplexity explicitly shows which memories contributed alongside
traditional web sources, making memory's impact auditable.  The
underlying architecture retrieves relevant context from a user's memory
store rather than injecting all memories, placing Perplexity in the
Advanced RAG category---the only consumer platform to exceed Naive RAG
for memory retrieval.  Like its peers, Perplexity lacks graph structure,
consolidation, and principled forgetting.

\paragraph{memU (NevaMind).}
memU~\cite{memu2026} is an agentic memory framework designed for 24/7
proactive agents.  Its distinguishing contribution is a hierarchical
file-system metaphor: raw interactions are transformed into structured,
searchable intelligence organized as a virtual file system, supporting
both embedding-based (RAG) and non-embedding (LLM) retrieval.  memU
receives multimodal inputs (conversations, documents, images) and
organizes them into daily notes, long-term memory, and extracted user
profiles.  A fourth ``intention layer'' (in development) analyzes
accumulated memories to predict user needs before they are expressed,
representing the most ambitious proactive memory capability in the
surveyed systems.  The framework claims approximately 10$\times$ token
reduction versus comparable usage patterns.  memU Bot extends the
framework into an enterprise-ready OpenClaw alternative with persistent
cross-session memory and team-scale deployment.

\paragraph{ReMe (MemoryScope).}
ReMe~\cite{reme2025} (formerly MemoryScope, ModelScope) provides a dual
memory architecture distinguishing \emph{personal memory} (user
preferences and traits) from \emph{task memory} (learned strategies from
prior executions).  Task memory implements four patterns: success
recognition, failure analysis, comparative learning across sampling
trajectories, and validation of extracted memories.  The framework
supports file-based and vector-based storage backends (Elasticsearch,
ChromaDB) with hybrid retrieval (vector + BM25).  ReMe integrates into
the AgentScope runtime as its default long-term memory service.  The
dual personal--task decomposition is architecturally distinct from most
systems' episodic--semantic split, treating memory \emph{function} rather
than memory \emph{type} as the primary organizing principle.

\paragraph{MIRIX.}
MIRIX~\cite{wang2025mirix} is a modular, multi-agent memory system with
six distinct memory types: Core, Episodic, Semantic, Procedural, Resource
Memory, and Knowledge Vault.  Its distinguishing contribution is
multimodal memory: unlike prior approaches confined to text, MIRIX
captures rich visual and multimodal experiences with temporally grounded
interactions.  A multi-agent framework dynamically coordinates updates
and retrieval across memory types.  On ScreenshotVQA, MIRIX achieves
35\% higher accuracy than the RAG baseline while reducing storage by
99.9\%; on LoCoMo, it attains 85.4\%, surpassing existing baselines.

%% ===============================================================
\subsection{Standalone Memory Servers}
\label{sec:standalone}

Several standalone servers provide memory as a network service,
typically via REST or MCP endpoints, without being tied to a specific
agent framework. \textbf{Motorhead} (Metal, Rust) implements a
minimal conversation buffer with incremental summarization on Redis:
when a configurable window (default 12 messages) fills, the oldest
half is summarized and the summary is incrementally updated.
\textbf{Engram} (Go) takes an agent-trusting approach---the LLM
client decides what to remember, storing curated summaries in SQLite
with FTS5 full-text search and no vector retrieval. Both represent
Naive RAG with no graph, no forgetting model, and no consolidation.
\textbf{OpenViking} (ByteDance/Volcengine, open-sourced January 2026)
adopts a virtual filesystem metaphor, mapping all agent context to a
\texttt{viking://} protocol with three tiers: L0 (immediate context),
L1 (session-level), and L2 (persistent archival). Built on VikingDB
infrastructure, OpenViking replaces fragmented vector storage with
hierarchical directory-based retrieval, but provides no forgetting,
no graph dynamics, and no preference learning.

\paragraph{MCP Memory Service.}
MCP Memory Service~\cite{doobidoo2026mcp} provides persistent memory
for AI agent pipelines (LangGraph, CrewAI, AutoGen, Claude) via a REST
API with knowledge graph and autonomous consolidation.  Storage uses
SQLite-vec with lightweight ONNX embeddings, achieving 5ms retrieval
latency.  The system supports over 20 AI client applications, making it
the most broadly compatible standalone memory server surveyed.  Its
\emph{Natural Memory Triggers} use AI-based heuristics (85\%+ accuracy)
combined with rule-based Context-Provider patterns (100\% guaranteed
coverage) to determine when to store new memories---a dual-strategy
approach unique among the surveyed systems.

\paragraph{Supermemory.}
\textbf{Supermemory}~\cite{supermemory2026} provides memory-as-a-service
via a REST API designed for multi-agent deployments.  The system
auto-learns from conversations, extracts structured facts, builds
evolving user profiles, handles contradictions through fact comparison,
and forgets expired information via TTL-based decay.  Its distinguishing
contribution is the ASMR (Adaptive Semantic Memory Retrieval) ensemble:
a single-pass retrieval achieves approximately 85.86\% on LongMemEval,
while an 8-variant ASMR ensemble reaches 99\%---among the highest
published scores on this benchmark.  Supermemory is a hosted service
with no self-hosted option, placing it in the Advanced RAG category.

\paragraph{Cognee.}
\textbf{Cognee}~\cite{cognee2026} (Topoteretes, Python) takes a
pipeline-based approach to memory construction: data is ingested from
30+ connectors (databases, APIs, file systems), processed through entity
and relationship extraction, and assembled into a knowledge graph
augmented with vector embeddings.  The system supports multimodal inputs
including images and audio.  Unlike most memory systems that operate
solely on conversation transcripts, Cognee's connector architecture
enables memory formation from diverse enterprise data sources.
Cognee is open source and provides both graph traversal and vector
similarity retrieval, placing it in the Advanced RAG category.  It
implements no forgetting or consolidation mechanisms.

\paragraph{EverMemOS.}
\textbf{EverMemOS}~\cite{evermemos2026} (EverMind AI, Python) implements
a self-organizing memory operating system with an engram-inspired
lifecycle: memories are created, strengthened, weakened, and
consolidated through use patterns analogous to biological engram
dynamics.  The architecture uses a dual-layer design:
\emph{working memory} performs real-time abstraction and context
management during active interactions, while \emph{long-term memory}
maintains dynamic knowledge graphs that restructure as new information
arrives.  EverMemOS is one of the few systems where the graph topology
evolves autonomously through use rather than being statically defined at
write time.  The system targets long-horizon reasoning tasks requiring
coherent behavior across extended interaction sequences.

\paragraph{MemOS.}
\textbf{MemOS}~\cite{memos2026} (MemTensor, Python) introduces the
concept of an AI memory operating system enabling persistent
\emph{skill memory} for cross-task reuse.  Its core abstraction is the
\emph{memory cube}---a composable, isolated memory unit that can be
shared across users, projects, and agents with controlled access
policies.  This isolation-and-sharing model is architecturally distinct
from the flat or tiered memory stores used by most systems: memory cubes
can be composed, versioned, and selectively exposed, enabling
fine-grained multi-tenant deployments.  MemOS is open source and targets
production scenarios where multiple agents must share learned knowledge
without cross-contamination.

\paragraph{OMEGA.}
\textbf{OMEGA}~\cite{omega2026} (omega-memory, Python) is a local-first
memory system for coding agents designed for zero-cloud operation: all
data, embeddings, and retrieval execute on the user's machine with no
external API calls required.  Storage uses a single SQLite database in
WAL mode extended with sqlite-vec for approximate nearest-neighbor search
and FTS5 for full-text keyword matching; embeddings are computed locally
via ONNX Runtime (\texttt{bge-small-en-v1.5} on CPU), eliminating
latency and privacy exposure from hosted embedding APIs.  The system
exposes 25 MCP tools covering working, episodic, and semantic memory
types, with automatic hook-based capture of decisions, error patterns,
and lessons during agent sessions.  Retrieval combines semantic vector
similarity, FTS5 keyword matching, and contextual priority scoring
with a reported 50\,ms latency.  A dynamic decay-curve forgetting
mechanism---rather than fixed TTLs---ages memories over time, and a
compression pass consolidates redundant facts.  OMEGA integrates
natively with Claude Code, Cursor, Windsurf, and Zed.  The developer
reports 95.4\% accuracy on LongMemEval, among the highest published
scores on that benchmark~\cite{omega2026}.  The design's primary
constraint is complete data locality---no telemetry, no cloud
calls, fully offline---differentiating it from hosted systems such
as Mem0 and Zep.  The system does not implement preference learning,
graph-overlay retrieval, or consolidation beyond compression.

%% ===============================================================
\subsection{Forgetting-Aware Systems}
\label{sec:forgetting-systems}

A cluster of recent systems (2025--2026) places principled forgetting at
the center of their architecture rather than treating it as an optional
lifecycle operation.  These systems draw explicitly from cognitive
science---particularly Ebbinghaus forgetting
curves~\cite{ebbinghaus1885} and retrieval-induced
forgetting~\cite{anderson1994}---to manage memory decay.

\paragraph{FadeMem.}
\textbf{FadeMem}~\cite{fademem2026} (Wei et al., Alibaba \& Peking
University) implements biologically inspired dual-layer forgetting via
adaptive exponential decay functions modulated by three signals: semantic
relevance, access frequency, and temporal patterns.  The architecture
separates a \emph{Short-term Memory Layer} (SML) for transient context
that fades quickly from a \emph{Long-term Memory Layer} (LML) for
high-importance facts such as user preferences and long-standing goals
that decay slowly.  An LLM-guided conflict resolution module detects
contradictions between incoming facts and existing memories, and a memory
fusion stage consolidates redundant entries while enforcing temporal
consistency.  Experiments span three benchmarks: on LTI-Bench (conflict
resolution), FadeMem achieves 68.9\% macro-averaged accuracy and 80.4\%
macro-averaged consistency; on LoCoMo, it attains a multi-hop F1 of
29.43 (vs.\ Mem0's 28.37 and MemGPT's 9.46); and after 30 days of
continuous interaction on the synthetic LTI-Bench dataset, it retains
82.1\% of critical facts while using only 55\% of the storage---compared
to Mem0's 78.4\% retention at 100\% storage.  This 45\% storage reduction
with superior retention represents the most favorable
storage-vs-retention tradeoff among forgetting-aware systems.  FadeMem
uses in-memory storage with no external database dependencies, making it
lightweight but unsuitable for persistent multi-session deployments
without additional infrastructure.  It does not implement graph
structure, preference learning, or structured retrieval beyond vector
similarity, placing it in the Naive RAG category despite its
sophisticated lifecycle management.

\paragraph{Oblivion.}
\textbf{Oblivion}~\cite{oblivion2026} (Rana, Hung et al., NEC
Laboratories Europe \& Georg-August-Universit\"at G\"ottingen) reframes
long-horizon agent memory as a \emph{control problem} rather than a
storage problem, decoupling memory into separate read and write paths.
The \emph{read path} uses uncertainty-gating: a Decayer computes
Ebbinghaus-inspired retention scores for each memory trace, and an
Activator triggers targeted retrieval only when the agent's uncertainty
exceeds a threshold and the working memory buffer is insufficient---
avoiding the ``always-on'' retrieval that plagues flat memory systems.
The \emph{write path} selectively reinforces memories that contributed to
forming the response, strengthening useful associations while allowing
unreinforced traces to fade in accessibility without explicit deletion.
The architecture separates transient \emph{Working Memory} from a
\emph{Persistent Store} organized hierarchically, maintaining high-level
strategies persistently while dynamically loading details as needed.
Oblivion is evaluated on LongMemEval (${\approx}$100K-token contexts)
and GoodAILTM (${\approx}$120K-token dynamic interactions) across four
LLM backbones (Phi-4-mini-instruct-4B, Qwen3-30B-A3B, GPT-4o-mini,
GPT-4.1-mini), consistently outperforming both direct and
memory-augmented baselines on both static and dynamic benchmarks.  The
read-path uncertainty gate is architecturally unique among the surveyed
systems.  However, Oblivion uses the same LLM for memory extraction,
uncertainty estimation, and response generation; smaller models with
weaker capability degrade retrieval-dependent tasks, and the framework
has not been evaluated on multilingual, multimodal, or domain-specific
settings.  Oblivion implements no graph structure, no consolidation
beyond reinforcement, and no preference learning.

\paragraph{MaRS.}
\textbf{MaRS}~\cite{mars2026} (Alqithami, Albaha University) introduces
privacy-aware memory management with six theoretically grounded
forgetting policies: FIFO, LRU, Priority Decay, Reflection-Summary,
Random-Drop, and a Hybrid variant.  The underlying \emph{Memory-Aware
Retention Schema} organizes four typed memory stores---episodic,
semantic, social, and task---as provenance-tracked nodes in a
graph-structured store, where edges encode temporal, causal, semantic,
and social relations.  Each node carries content, timestamps, token-cost
weights, sensitivity scores, and provenance fields recording how the item
entered memory, enabling agents to justify retention or deletion relative
to source reliability and consent status.  Retention is formulated as a
resource-allocation problem under explicit token budgets: forgetting
policies jointly consider access statistics, semantic centrality, task
relevance, and privacy sensitivity.  Each policy is analyzed for
computational complexity and sensitivity-aware retention, with optional
$(\varepsilon, \delta)$-differential privacy guarantees via calibrated
noise injection.  The accompanying \emph{FiFA benchmark} (Forgetful but
Faithful Agent) evaluates 300 runs across multiple memory budgets and
agent configurations on five dimensions: narrative coherence, goal
completion, social recall accuracy, privacy preservation, and cost
efficiency.  The Hybrid policy achieves a 0.911 composite score,
consistently preserving coherence and social recall while keeping
leakage low and costs tractable.  MaRS is the only surveyed system that
formalizes forgetting as a privacy mechanism with provable guarantees
rather than treating it solely as a storage optimization.  Its
limitation is the absence of a vector retrieval component or embedding-
based similarity search---retrieval relies on the typed graph indices
and policy-driven filtering, which may underperform on open-ended
semantic queries compared to hybrid retrieval systems.

\paragraph{SuperLocalMemory V3.3.}
\textbf{SuperLocalMemory V3.3}~\cite{superlocalmemory2026} (Bhardwaj,
independent researcher) implements 7-channel cognitive retrieval
combining semantic, keyword, entity graph, temporal, spreading
activation, consolidation, and Hopfield associative channels with
Ebbinghaus-inspired adaptive forgetting curves coupled to progressive
embedding compression.  The system builds on the information-geometric
foundations of V3.0: Fisher-Rao geodesic
distance for quantization-aware similarity (via the novel FRQAD metric,
achieving 100\% precision at preferring high-fidelity embeddings over
quantized ones, vs.\ 85.6\% for cosine similarity), cellular sheaf
cohomology for contradiction detection, and Riemannian Langevin dynamics
for stochastic lifecycle management.  V3.3 adds \emph{long-term implicit
memory} via soft prompts that parameterize learned patterns, and a
zero-friction auto-cognitive pipeline that implements the complete memory
lifecycle---recall, observe, save, consolidate, forget,
parameterize---through seven Claude Code hooks with zero configuration.
The system's defining constraint is \emph{zero-LLM operation} (Mode~A):
all memory management, retrieval, and forgetting run without any LLM
calls, using only local models (SQLite with sqlite-vec, local ONNX
embeddings) and classical algorithms.  This achieves 70.4\% on LoCoMo
without cloud dependencies---the highest score among systems that operate
entirely without LLM involvement.  An optional Mode~C with LLM
assistance reaches 87.7\%.  The system ships as a single
\texttt{npm install -g superlocalmemory} package compatible with 17+
AI tools (Claude Code, Cursor, Windsurf, Copilot, Gemini CLI, etc.).
The zero-LLM constraint makes SuperLocalMemory suitable for
privacy-sensitive and offline deployments, but the system lacks
multi-user support, has no graph dynamics beyond the static entity graph,
and the mathematical sophistication (information geometry, sheaf theory)
may impede community adoption and extension.

%% ===============================================================
\subsection{Framework-Level Memory}
\label{sec:framework}

The major agent frameworks provide memory as a built-in module, but
these integrations are typically shallow---memory as an afterthought
rather than a first-class architectural concern. \textbf{LangChain}'s
memory classes (ConversationBufferMemory, ConversationSummaryMemory,
ConversationBufferWindowMemory) offer working-memory-only storage with
truncation-based forgetting; most are now deprecated in favor of
\texttt{RunnableWithMessageHistory}. \textbf{LlamaIndex} provides
composable memory blocks (StaticMemoryBlock,
FactExtractionMemoryBlock, VectorMemoryBlock) with FIFO eviction and
a configurable token ratio (70\% chat history, 30\% long-term by
default). \textbf{LangMem SDK}, LangChain's dedicated memory engine,
adds semantic memory (extracted facts and preferences) and a
distinctive \emph{procedural} memory capability: learned patterns are
saved as updated prompt instructions, modifying the agent's behavior
without explicit retrieval. None of these frameworks implements a
graph overlay, principled forgetting, consolidation, or preference
emergence. Their retrieval ranges from direct prompt injection
(LangChain) to vector similarity via LangGraph's storage layer
(LangMem), placing them uniformly in the Naive RAG category.

Two frameworks stand apart with substantially richer memory modules.
\textbf{AgentScope}~\cite{agentscope2025} (Alibaba/ModelScope) provides
layered memory: in-memory short-term storage, Redis-backed working memory
with TTL-based expiration, LLM-driven memory compression (added January
2026), and ReMe-based long-term memory with hybrid vector+BM25 retrieval.
AgentScope is one of the few frameworks with both consolidation and
forgetting, placing it in the Advanced RAG category.
\textbf{CAMEL}~\cite{camel2025} implements composable memory blocks via
the Composite pattern: \texttt{ChatHistoryMemory} (windowed history with
configurable window size), \texttt{VectorDBMemory} (vector similarity
over past interactions), and \texttt{LongtermAgentMemory} (a combined
module that augments chat history with vector retrieval).  CAMEL provides
the most flexible memory composition in the framework category, though
it lacks graph structure, consolidation, and forgetting.

%% ===============================================================
\subsection{Coding Agent Memory}
\label{sec:coding-agents}

A distinctive cluster of memory systems has emerged around coding
agents (Claude Code, Cursor, Windsurf, Augment Code), exposed primarily
via the Model Context Protocol (MCP).  While these tools originated as
coding assistants, their deployment as general-purpose agents---
particularly through MCP-based memory extensions---means their memory
architectures must serve domain-agnostic use cases.  We retain the
``coding agent'' label to reflect their primary deployment context, while
noting that the memory systems built for them (memsearch, QMD, Basic
Memory) are architecturally indistinguishable from dedicated memory
systems.  These systems share a common pattern:
replace the default full-context injection of MEMORY.md files with
ranked retrieval, typically achieving 70--96\% token savings.
\textbf{memsearch}~\cite{zilliz2026a} (Zilliz/Milvus team) extracts
the default memory file structure and overlays BM25 plus vector
retrieval with RRF fusion and progressive top-$k$ disclosure.
\textbf{QMD}~\cite{lutke2026qmd} (Tobi L\"utke, Shopify) runs three
local GGUF models for embedding (Gemma-300M), reranking
(Qwen3-Reranker-0.6B), and query expansion (1.7B parameter model),
achieving fully offline hybrid retrieval with cross-encoder reranking
and reported 96\% token reduction~\cite{levine2026}. \textbf{Beads}
(Go, on Dolt version-controlled SQL) models a dependency-aware issue
graph with four link types, enabling traversal along dependency chains
rather than pure similarity. \textbf{Basic Memory} (Python) takes a
wiki-link approach, storing structured Markdown with explicit
Observations and Relations that form a semantic knowledge graph
navigable via \texttt{memory://} URLs. All four independently adopted
RRF over the default weighted linear fusion---a convergence we
analyze further in \Cref{sec:memorymd}.

The OpenClaw ecosystem has spawned multiple alternative agents, each with
distinct memory approaches.
\textbf{OpenFang}~\cite{openfang2026} (RightNow AI, Rust) is the most
architecturally ambitious: its dedicated \texttt{openfang-memory} crate
implements SQLite-backed storage with vector embeddings, a knowledge
graph, cross-channel canonical sessions (context persists across
Telegram, Slack, Discord, etc.), and automatic LLM-based compaction.
OpenFang ships as a single 32MB binary with 40 channel adapters, making
it the most infrastructure-light system with graph capabilities.
\textbf{IronClaw}~\cite{ironclaw2026} (NEAR AI, Rust) implements
PostgreSQL with pgvector for persistent vector search, Reciprocal Rank
Fusion combining full-text search (tsvector/tsquery) and vector
similarity (cosine distance), and WASM sandboxing for all untrusted
tools.  Its credential isolation architecture (secrets are never exposed
to the LLM, injected only at the host boundary) is unique among the
surveyed systems.
\textbf{Nanobot}~\cite{nanobot2026} (HKU, Python) takes the opposite
approach: approximately 4,000 lines of core code with a
\texttt{MemoryStore} class maintaining daily notes
(\texttt{memory/YYYY-MM-DD.md}) and long-term memory
(\texttt{MEMORY.md})---no vector retrieval, no graph, just two plain
files and grep.

Four IDE-integrated memory systems emerged in late 2025 and early 2026.
\textbf{GitHub Copilot Memory}~\cite{copilot2026} introduces a
\emph{citation-based} architecture: memories are stored with references
to specific code locations, and agents verify citations in real-time
before using a memory.  In adversarial testing, agents consistently
detected contradictions between memories and code, updated incorrect
memories, and the memory pool self-healed.  This is the only system in
our survey that implements contradiction resolution through citation
verification rather than LLM-based fact comparison.  Memories
automatically expire after 28 days and are refreshed only through
re-validation.
\textbf{Cursor} provides two memory mechanisms:
\emph{Cursor Memory}~\cite{cursor2026} (IDE-level RAG indexing of the
entire repository via \texttt{@Codebase} symbols with vector embeddings)
and \emph{Cursor Rules} (persistent project/user/team instructions in
\texttt{.cursor/rules/*.mdc} files, injected as file-based context).
\textbf{Windsurf Memory}~\cite{windsurf2026} (Cascade) auto-generates
memories during conversation when it encounters useful context, with
a user toggle for autonomous memory creation and
\texttt{.windsurfrules} for persistent project rules.
\textbf{Augment Code}~\cite{augment2026} is the most
infrastructure-heavy IDE memory system, designed for enterprise
codebases exceeding 400{,}000 files.  Its Context Engine performs
real-time repository indexing across multiple repositories, tracking
cross-service dependencies and commit history.  The \emph{Memories}
feature stores project-specific facts---architecture decisions, coding
patterns, preferred libraries---that persist across sessions and are
automatically created by the agent when it encounters reusable context.
A \emph{Memory Review} interface lets users approve, edit, or reject
memories before finalization, providing a manual curation loop unique
among IDE memory systems.  Despite this sophistication, Augment's
memories lack automated lifecycle management: there is no consolidation,
no forgetting mechanism, and no contradiction resolution---memories
persist indefinitely until manually deleted.

\textbf{episodic-memory}~\cite{vincent2025episodic} (Jesse Vincent,
Superpowers ecosystem) archives Claude Code's conversation JSONL logs
into a SQLite database with \texttt{sqlite-vec} vector embeddings
(generated locally via Transformers.js), enabling semantic search over
the complete interaction history.  Its distinguishing contribution is a
\emph{subagent delegation} pattern: a specialized Haiku subagent handles
all memory retrieval and context assembly, preventing the context bloat
that arises from injecting raw conversation transcripts into the primary
agent's context window.  The system implements a session-end learning
loop in which the agent reflects on completed work---a ``reflect''
pattern independently adopted by several other systems.  The plugin is
part of the Superpowers marketplace, one of the largest Claude Code
plugin ecosystems.

\textbf{claude-mem}~\cite{newman2025claudemem} captures tool usage and
agent actions across five lifecycle hooks
(SessionStart, UserPromptSubmit, PostToolUse, Summary, SessionEnd) and
compresses raw observations into structured knowledge using Claude's
Agent SDK---a form of automated episodic-to-semantic consolidation.
Retrieval follows a token-efficient three-layer workflow:
\emph{search} returns a compact index of matching observation IDs,
\emph{review} presents titles and metadata for filtering, and
\emph{fetch} retrieves full details only for selected observations.
This progressive disclosure pattern minimizes context consumption---the
agent reads observation titles ($\sim$100 tokens) before deciding
whether to load full content ($\sim$500 tokens each).  An experimental
\emph{Endless Mode} implements a biomimetic memory architecture for
extended sessions, automatically managing context pressure through
adaptive summarization.  With over 28{,}000 GitHub stars, claude-mem is
among the most widely deployed agent memory systems.

\textbf{MemPalace}~\cite{mempalace2026} (Jovovich \& Sigman, Python)
challenges the consolidation-first consensus with a deliberately
anti-extractive architecture: raw verbatim transcripts are stored in
ChromaDB without summarization, and a hierarchical \emph{palace}
structure---wings (projects), halls (memory types), rooms
(topics)---scopes retrieval before embedding similarity is applied.
This spatial narrowing yields a 34\% retrieval improvement over flat
search on 22{,}000+ real memories.  On LongMemEval the raw mode scores
96.6\% R@5 without any LLM involvement, and a hybrid mode with Haiku
reranking reaches 100\% (500/500)---the highest published score on this
benchmark.  An experimental compression dialect (AAAK, $\sim$30$\times$
token reduction) is under development but currently regresses retrieval
to 84.2\%.  MemPalace implements no consolidation, forgetting, or
contradiction resolution, making it the purest test of the hypothesis
that raw episodic storage with structured retrieval can outperform
systems that invest in semantic extraction.  With over 27{,}000 GitHub
stars within days of release, it has become one of the most visible
entrants in the space.

\textbf{AgentMemory}~\cite{agentmemory2026} (Jordan McCann, Python)
implements a complete memory operating system with a knowledge graph,
a multi-stage consolidation pipeline, contradiction detection, and
confidence decay.  The storage backend uses SQLite by default (with
PostgreSQL available for multi-process deployments), and the retrieval
engine combines six weighted signals via a custom fusion pipeline:
semantic similarity (0.30), lexical/BM25 (0.12), spreading activation
(0.18), graph traversal (0.18), importance scoring (0.10), and temporal
grounding (0.12).  A deterministic HNSW index (using SHA-256 vector
hashing for content-based level assignment) provides reproducible nearest
neighbor search, and a cross-encoder reranker refines the final ranking.
The consolidation pipeline is the most comprehensive among coding agent
memory systems: memories are progressively abstracted from raw
observations through episodic summaries to semantic knowledge, with
near-duplicate merging and contradiction detection at each promotion
stage.  AgentMemory achieves 96.2\% (481/500) on LongMemEval in a single
deterministic run (\texttt{PYTHONHASHSEED=42}, no ensembling, no oracle
access), placing it among the highest single-pass scores on this
benchmark---surpassing Mastra (93.27\%), Supermemory single-pass
(85.86\%), and Zep.  The system exposes both a FastAPI REST server and an
MCP server for native tool use in Claude Desktop and compatible hosts.
Built by a solo developer in 16 days on a mid-range gaming PC (Intel
Core i3-12100F) for approximately \$1{,}000 in API costs, it
demonstrates that a single-developer effort can match or exceed the
benchmark performance of team-built production systems.  AgentMemory does
not implement principled forgetting beyond confidence decay, and its
six-signal retrieval requires careful weight tuning that may not
generalize across domains.

\textbf{Capy-Cortex}~\cite{capycortex2026} (Python) provides autonomous
memory for coding agents with a triple fusion retrieval pipeline
combining FTS5 full-text search, TF-IDF embeddings, and entity graph
traversal, merged via Reciprocal Rank Fusion (RRF) with sub-10ms
retrieval latency.  Storage uses SQLite with FTS5, and the entity
knowledge graph tracks relationships between tools, libraries, errors,
and concepts with temporal version chains recording how knowledge
evolves.  The system's distinguishing feature is \emph{quality-gated
storage}: every extracted insight is scored 0--4 on four dimensions
(actionable, specific, novel, durable) by a fast LLM, with a threshold
of 2/4---anything below is discarded, preventing the noise accumulation
that plagues most learning systems.  A dual-model architecture uses a
capable model (e.g., GPT-4o, Claude Sonnet) for causal analysis of
failures and a fast model (e.g., GPT-4o-mini, Haiku) for quality
scoring.  Seven lifecycle hooks capture agent events automatically: on
tool failure, the smart model performs root cause analysis and generates
a rule; on tool success, helpful rules receive a confidence boost; at
session end, learnings are extracted and stored.  A self-curating
consolidation process (every 10 sessions) clusters related memories,
merges near-duplicates (cosine similarity threshold 0.85), and
synthesizes abstract principles, while unreinforced rules undergo
confidence decay with a 90-day half-life.  The combination of quality
gating at write time, success-based reinforcement during interaction, and
periodic consolidation makes Capy-Cortex one of the few coding agent
memory systems with both ingestion filtering and lifecycle management.
The system does not publish benchmark results on standard evaluations
(LoCoMo, LongMemEval), and its effectiveness depends heavily on the
quality of the underlying LLM pair for causal analysis and scoring.

\paragraph{chum-mem.}
\textbf{chum-mem}~\cite{chummem2026} (sly-codechum, TypeScript/Rust)
implements a provenance-aware memory platform for coding agents
structured around the \emph{Proof-Carrying Knowledge Compiler} (PCKC)
model.  Sessions are segmented into episodes (conversation,
implementation, debugging), from which atomic \emph{claims} of eight
types---fact, decision, task, constraint, bug, fix, implementation
detail, open question---are extracted with \emph{authority
classifications}: tool-verified, user-confirmed, repository-derived,
session-derived, or model-inferred.  Each claim is attached to a
structured \emph{proof}---a verifiable chain of evidence linking the
assertion to its source events---enabling the system to distinguish
what was directly observed from what was inferred.  A \emph{belief gate}
hard-rejects model-generated prose and unverified agent reasoning at
write time: only claims carrying a non-model-inferred authority class
enter the knowledge store.  A contradiction engine and supersession
engine manage conflicts: when two claims disagree, the newer supersedes
the older and both are retained with explicit \texttt{superseded} status,
preserving full reasoning history.  Leiden community detection organizes
the resulting graph into coherent topic clusters.  Retrieval scores
candidates on eight signals (lexical, semantic, session relevance, graph
proximity, recency, importance, confidence, freshness/supersession
penalties) and assembles a \emph{compiled context}: a token-budgeted
minimal proof set rather than a top-$k$ similarity list.  The backend
uses Rust/Axum with PostgreSQL (pgvector, FTS, row-level security for
multi-tenancy) and optional ChromaDB for typed embedding partitions.
chum-mem is architecturally notable for treating memory correctness as a
formal property: claims require proof provenance before admission, and
query answers are verified deductions from stored evidence rather than
plausibility rankings.  The system does not yet publish results on
standard benchmarks (LoCoMo, LongMemEval).

\paragraph{bellamem.}
\textbf{bellamem}~\cite{bellamem2026} (immartian, TypeScript v0.3)
implements a persistent \emph{belief-graph forest} for coding agents,
grounding memory in Bayesian mass accumulation rather than similarity
search.  Every concept accumulates \emph{R1 mass} via a frozen
Jaynes log-odds formula: the first citation from a new voice contributes
$+0.5$ log-odds; repeat citations from the same speaker contribute only
$+0.1$, rewarding consensus over repetition.  Mass values are
$\sigma$-squashed into $[10^{-6},\,1-10^{-6}]$ to prevent saturation,
and every delta is traceable to the specific turn that caused it.
Concepts carry a dual-axis classification: \emph{class}---invariant,
decision, observation, ephemeral---captures temporal profile;
\emph{nature}---factual, normative, metaphysical---captures epistemic
type.  The two axes enable precise retrieval: an invariant-factual
entry (``Python 3.11 removed module X'') and a decision-normative entry
(``we decided to skip Python 3.11'') are indexed and filtered
independently.  The graph supports seven typed edge kinds; crucially,
\emph{dispute} and \emph{retract} edges are first-class.  A dispute
edge is created when two sources contradict each other on the same
concept---it is preserved permanently rather than resolved, recording
the conflict as a durable artifact.  A retract edge is created when a
speaker withdraws a prior claim, marking that concept
\texttt{retracted} without erasing the record.  Retrieval executes an
\emph{expand} operation: seed concepts are scored by
$\max(\text{substring match, cosine}) \times \text{mass}$, then a
1-hop edge walk collects neighbors grouped by edge type---disputes
first, then causes, support, elaborate, retract---returning structural
explanations rather than isolated similarity hits.  An R5 sweep ages
open ephemerals to \texttt{stale} after seven days of inactivity.
bellamem is architecturally notable for treating disagreement as
information: contested beliefs and withdrawn claims are preserved and
surfaced during retrieval, enabling agents to reason about the
provenance and confidence of their own knowledge.  No benchmark results
on standard evaluations have been published.

\paragraph{Memori.}
\textbf{Memori}~\cite{memori2026} (Memori Labs, multi-framework) adopts
a \emph{memory-from-execution} paradigm: rather than extracting
knowledge from conversation text alone, the system captures agent tool
calls, decisions, outcomes, and execution state transitions as
first-class memory artifacts.  The storage backend uses a SQL-native
relational schema (Postgres or MySQL) with normalized tables for facts,
entities, events, preferences, and policies---replacing unstructured
text blobs with explicit typed columns, foreign-key relationships, and
auditable provenance.  Retrieval reconstructs context by joining these
normalized tables with schema-level constraints rather than semantic
similarity, which the authors report reduces hallucination risk and
inference cost: Memori achieves 81.95\% on LoCoMo using approximately
1{,}294 tokens per query, approximately 5\% of the full-context token
budget.  Multi-agent coordination handles read-time conflict management
and iterative revision.  Memori's execution-trace paradigm is
architecturally complementary to episode-level systems: where episodic
memory records \emph{what was said}, Memori records \emph{what was done
and what happened}---tool call results, state changes, and observable
outcomes with explicit timestamps.  The system does not implement
forgetting, graph traversal, or preference learning beyond execution
traces.

Despite their retrieval sophistication, no IDE-integrated memory system
implements consolidation or automated forgetting.  Copilot's 28-day
expiry is the closest approximation, but it is a fixed TTL rather than
usage-adaptive decay.  This lifecycle gap will become acute as memory
stores grow unboundedly---particularly for Augment, where enterprise
teams accumulating memories across hundreds of thousands of files face
eventual retrieval degradation without principled memory management.

%% ===============================================================
\subsection{File-Based Memory}
\label{sec:filebased}

The simplest memory architecture stores knowledge in Markdown files
that the agent reads and writes directly. The dominant instance is the
\textbf{MEMORY.md pattern} used by coding agents such as OpenClaw:
a curated long-term knowledge file (MEMORY.md) plus daily interaction
logs with temporal decay, backed by a SQLite index with FTS5 and
sqlite-vec for hybrid retrieval (70\% vector, 30\% BM25 weighted
fusion). \textbf{Claudesidian} extends this pattern into Obsidian
vaults, using PARA folder organization and daily/weekly notes with
Obsidian's wiki-link graph for navigation. File-based memory provides
zero-infrastructure deployment and human-readable persistence, but
scales poorly: MEMORY.md is injected into the system prompt in full
on every turn regardless of relevance, consuming approximately 35,600
tokens per message (93.5\% waste in multi-message
sessions)~\cite{openclaw9157}. We analyze the cost, structure, and
retrieval failure modes of file-based memory in detail in
\Cref{sec:memorymd}.

%% ===============================================================
\subsection{Research Architectures}
\label{sec:research}

\paragraph{Generative Agents.}
Park et al.~\cite{park2023generative} introduced the seminal memory
architecture for LLM agents: a chronological \emph{memory stream} of
observations, from which the agent periodically generates
\emph{reflections}---higher-order thoughts that synthesize multiple
observations into abstract insights. Retrieval scores each memory by a
weighted product of three signals: \emph{recency} (exponential decay),
\emph{importance} (LLM-rated 1--10), and \emph{relevance} (cosine
similarity to the current query). This triple-scored formula became the
foundational retrieval mechanism for agent memory, adopted or adapted
by the majority of subsequent systems.

The architecture's significance lies less in its individual components
(all relatively simple) than in their composition: the reflection
mechanism implements a form of episodic-to-semantic consolidation,
where accumulated observations are distilled into reusable knowledge.
The system uses in-memory storage (ephemeral), a flat memory stream
(no typed stores), no graph overlay, and single-curve recency
decay---limitations that subsequent work has systematically addressed.
The reflection mechanism, however, inspired consolidation pipelines
in multiple production and research systems.

\paragraph{HippoRAG.}
HippoRAG~\cite{gutierrez2024hipporag} grounds its architecture in
hippocampal indexing theory from neuroscience: the LLM acts as
neocortex (processing and generating language), a parahippocampal
region (PHR) encoder detects synonymy and entity co-reference, and an
open knowledge graph acts as hippocampus (indexing and associating
memories). Retrieval uses Personalized PageRank (PPR) on the knowledge
graph, propagating activation from query-matched nodes to
transitively connected entities. This enables multi-hop reasoning that
pure embedding similarity cannot achieve---asking ``who handles auth
permissions?'' can retrieve information about ``Alice manages the auth
team'' through entity-relationship traversal rather than lexical
overlap.

HippoRAG achieves 20\% improvement on multi-hop question answering at
10--30$\times$ lower cost than iterative retrieval
approaches. HippoRAG~2~\cite{gutierrez2025rag2memory} (ICML~2025)
extends this with deeper passage integration---full passages become
first-class graph nodes alongside phrase nodes---and more effective
online LLM use for knowledge graph construction, adding episodic
storage and hybrid fusion (PPR + dense embedding) to the original
semantic-plus-graph design.  Neither version implements forgetting,
consolidation, or preference learning; both are retrieval
architectures rather than full memory lifecycle systems. PPR on a
knowledge graph is comparable in cognitive inspiration to spreading
activation on associative graphs, though the two mechanisms differ:
PPR computes a global stationary distribution, while spreading
activation models local, time-bounded propagation~\cite{collins1975}.

\paragraph{SYNAPSE.}
SYNAPSE~\cite{jiang2026synapse} constructs a unified episodic-semantic
graph and retrieves through two biologically inspired mechanisms:
\emph{spreading activation} (activation propagates from query-matched
nodes to neighbors through weighted edges, priming associated
memories) and \emph{lateral inhibition} (strongly activated nodes
suppress weakly activated competitors, sharpening retrieval focus).
Retrieval fuses geometric embedding similarity with activation-based
graph traversal in a triple-hybrid pipeline. Temporal decay weakens
links over time.

SYNAPSE achieves +7.2 F1 over prior state of the art on LoCoMo, 23\%
improvement on multi-hop reasoning, and 95\% token reduction versus
full-context approaches. The lateral inhibition mechanism is
particularly notable: it implements a form of retrieval-induced
forgetting~\cite{anderson1994}, where retrieving one memory
suppresses competing memories sharing the same cues. This is the most
neurobiologically grounded retrieval mechanism in the surveyed
systems, though SYNAPSE does not implement consolidation, principled
forgetting (beyond temporal decay), or preference learning. The
system is in-memory Python research code, not a persistent or
embeddable library.

\paragraph{Mem-alpha.}
Mem-alpha~\cite{wang2025memalpha} takes a fundamentally different
approach to memory management: instead of hand-crafting lifecycle
heuristics, it trains memory policies via reinforcement learning.
The architecture decomposes memory into three components---core
(always in context), episodic (interaction history), and semantic
(distilled knowledge)---a decomposition nearly identical to the CoALA
framework. An RL agent learns when to store, what to consolidate,
and what to forget, using downstream question-answering accuracy as
the reward signal. Trained on 30K tokens, the learned policy
generalizes to 400K+ tokens (13$\times$ extrapolation).

The RL-learned approach represents the most significant methodological
divergence from hand-crafted cognitive models. Learned policies can
potentially discover non-obvious management strategies that human
designers would not specify, but they require training data, are
opaque (the policy network's decisions are not interpretable), and may
not transfer across domains. Mem-alpha is part of a broader 2025--2026
trend: Memory-R1, MemRL~\cite{zhang2026memrl}, and AgeMem all train
memory management via RL, collectively establishing learned memory
policies as a distinct paradigm alongside rule-based and
neuroscience-inspired approaches.

\paragraph{SimpleMem.}
SimpleMem~\cite{simplemem2026} introduces an efficient memory framework
based on semantic lossless compression, implemented as a three-stage
pipeline: (1)~Semantic Structured Compression, which distills
unstructured interactions into compact, multi-view indexed memory units;
(2)~Online Semantic Synthesis, an intra-session process that integrates
related context into unified abstract representations to eliminate
redundancy; and (3)~Intent-Aware Retrieval Planning, which infers search
intent to dynamically determine retrieval scope.  SimpleMem achieves a
26.4\% average F1 improvement over baselines while reducing
inference-time token consumption by up to 30$\times$.

\paragraph{AriadneMem.}
AriadneMem~\cite{ariadnemem2026} extends graph-based memory with two
innovations: \emph{conflict-aware coarsening} for memory updates that
detects and resolves contradictions during graph construction, and
\emph{approximate Steiner completion with bridge-node discovery} to
construct connected evidence subgraphs for retrieval.  Rather than
expanding to large contexts or running iterative planning loops,
AriadneMem retrieves a compact connected subgraph and exposes explicit
reasoning paths to the generator.  It achieves the highest average F1
across all evaluated backbones (42.57 on GPT-4o vs.\ 39.06 for SimpleMem
and 36.09 for Mem0), with the largest gains on multi-hop and temporal
reasoning subsets.

\paragraph{MemR3.}
MemR3~\cite{memr3_2025} introduces closed-loop control over memory
retrieval through two mechanisms: a router that selects among retrieve,
reflect, and answer actions, and a global evidence-gap tracker that
renders the answering process transparent.  This design departs from the
standard retrieve-then-answer pipeline, minimizing unnecessary lookups
while providing a human-readable trace of the agent's decision process.

\paragraph{PersonaRAG.}
PersonaRAG~\cite{personarag2024} incorporates user-centric agents to
adapt retrieval and generation based on real-time user data.  Four
specialized agents---engagement tracking, preference analysis, context
understanding, and feedback integration---analyze interactions from
different perspectives, achieving over 5\% accuracy improvement over
baseline RAG.

\paragraph{ID-RAG.}
ID-RAG~\cite{idrag2025} addresses identity drift in long-horizon
generative agents by grounding persona in a structured identity
knowledge graph of core beliefs, traits, and values.  During the agent's
decision loop, this graph is queried to retrieve relevant identity
context, directly informing action selection.  The separation of identity
memory from other memory types is architecturally distinctive and
addresses a failure mode (persona drift) that no other surveyed system
explicitly targets.

\paragraph{MAGMA.}
MAGMA~\cite{magma2026} (Jiang et al., UT Dallas \& University of
Florida) introduces a multi-graph architecture that represents each
memory item across four orthogonal graph views: semantic (conceptual
relationships), temporal (time-ordered event sequences with resolved
relative dates), causal (cause-effect chains), and entity (typed entities
and attributes).  An \emph{Adaptive Traversal Policy} formulates
retrieval as heuristic beam search over these relational views: the
policy assigns intent-specific weights to edge types (e.g., high weight
on causal edges for ``why'' queries, temporal edges for ``when''
queries), combining structural graph traversal scores with embedding
similarity at each step to retain the top-$k$ most relevant nodes.  A
\emph{dual-stream memory evolution} mechanism decouples latency-sensitive
event ingestion from asynchronous structural consolidation, preserving
responsiveness while refining relational structure in the background.
On LoCoMo, MAGMA achieves an overall LLM-as-Judge score of 0.70
(vs.\ 0.48 for A-MEM and 0.58 for Nemori), with the largest gains on
single-hop retrieval (0.551~F1) and temporal reasoning (0.509~F1).  On
LongMemEval, MAGMA reaches 61.2\% average accuracy while reducing token
consumption by 95\% relative to full-context approaches and achieving
query latency of 1.47s (40\% faster than the next best baseline).  The
system is implemented as in-memory Python research code with open-source
availability.  MAGMA's architectural insight---that different query types
benefit from different graph structures---is complementary to systems
like SYNAPSE (which uses spreading activation on a unified graph) and
HippoRAG (which uses Personalized PageRank on a single knowledge graph).
However, MAGMA does not implement forgetting, preference learning, or
lifecycle management beyond the asynchronous consolidation stream, and
its four-graph representation requires higher memory overhead than
single-graph approaches.

\paragraph{Mastra Observational Memory.}
Mastra's Observational Memory~\cite{mastra2026} replaces traditional
retrieval-augmented memory with a continuous compression architecture
using two background agents.  An \emph{Observer} agent monitors growing
message history and converts raw conversations into dense, timestamped
observations prioritized by a three-tier scheme (red~=~critical,
yellow~=~important, green~=~informational); the original messages
are then dropped and replaced with these observations.  When observations
themselves exceed a configurable threshold (40{,}000 tokens by default),
a \emph{Reflector} agent condenses them further, combining related
entries and dropping stale information---a two-stage compression pipeline
that achieves up to 40$\times$ token reduction.  Crucially, the main
agent never directly reads from or writes to memory; memory formation is
fully autonomous via the background agents.  The resulting context window
is predictable, low-latency, and fully prompt-cacheable---no per-turn
dynamic retrieval, no vector database, no graph store.  This text-only
architecture is compatible with prompt caching from Anthropic and OpenAI,
yielding up to 10$\times$ cost reduction in production deployments.  The
key finding is that observations outperform raw oracle data: retrieving
curated observations yields higher accuracy than retrieving the original
source material, suggesting that the extraction process itself adds value
through denoising and abstraction.  On LongMemEval, Mastra achieves
84.23\% with GPT-4o, 93.27\% with Gemini-3-pro-preview, and 94.87\%
with GPT-5-mini---the latter exceeding the oracle baseline.  The system
ships as part of Mastra~1.0 with plug-ins for LangChain, Vercel AI~SDK,
and other frameworks.  Mastra implements no forgetting, graph structure,
or consolidation beyond the Observer/Reflector pipeline, and its
effectiveness scales with the quality of the underlying LLM: the 10-point
gap between GPT-4o and GPT-5-mini scores demonstrates strong model
dependence.

\paragraph{GAM.}
\textbf{GAM} (Hierarchical Graph-based Agentic Memory)~\cite{gam2026}
(Wu et al.) introduces a two-tier hierarchical graph architecture that
explicitly decouples memory encoding from consolidation.  An
\emph{event progression graph} accumulates ongoing dialogue with low
latency, capturing transient incremental context.  A
\emph{topic associative network} receives triggered consolidation bursts:
when semantic drift is detected in the event progression graph, a
consolidation cycle distills accumulated episodes into persistent topic
nodes and associative edges.  This sleep-inspired cycle mirrors
hippocampal--neocortical consolidation: rapid encoding in a
high-plasticity episodic buffer followed by asynchronous integration
into a stable semantic store.  The decoupling resolves the fundamental
tension between rapid context perception---which stream-based systems
handle well but suffer from noise accumulation over time---and stable
long-term retention---which discrete consolidation architectures
provide but at the cost of narrative flexibility.  Retrieval applies a
graph-guided multi-factor scoring strategy that traverses both tiers,
weighting candidates by recency, semantic relevance, and inter-topic
association strength.  GAM outperforms baselines on LoCoMo and
LongDialQA in both reasoning accuracy and retrieval efficiency.  The
system is implemented as Python research code.

\paragraph{Kumiho.}
\textbf{Kumiho}~\cite{kumiho2026} (Park) presents the first agent memory
system grounded in formal \emph{AGM belief revision semantics}
(Alch\'{o}urron, G\"{a}rdenfors, and Makinson~\cite{agm1985}), proving
that its versioned graph architecture satisfies the K*2--K*6 revision
postulates and both Hansson Relevance and Core-Retainment postulates for
safe contraction.  The dual-store design pairs Redis working memory with
a Neo4j long-term graph.  Graph primitives are immutable revisions with
mutable tag pointers: nodes are never overwritten; instead, a new
revision is linked to its predecessor and tag pointers advance forward,
providing a complete temporal audit trail.  Typed epistemic dependency
edges capture causal, logical, and provenance relationships.  Three
retrieval enhancements are introduced: (1) \emph{prospective indexing}
generates LLM-predicted future-scenario implications at write time,
indexing anticipated consequences alongside recorded facts; (2)
\emph{event extraction} preserves structured causal events rather than
compressing them away during summarization; and (3)
\emph{client-side LLM reranking} re-scores retrieved candidates with
query context before final selection.  On LoCoMo, Kumiho achieves
0.565~F1 ($n=1{,}986$) with 97.5\% adversarial refusal; on LoCoMo-Plus,
it achieves 93.3\% judge accuracy ($n=401$) versus a Gemini~2.5 Pro
baseline of 45.7\%.  By grounding its versioning scheme in formal belief
revision theory, Kumiho provides correctness guarantees absent from
ad hoc contradiction-resolution strategies employed in systems such as
Mem0 and AgentMemory.

\paragraph{A-MEM.}
\textbf{A-MEM} (Agentic Memory)~\cite{amem2025} (Xu et al.,
NeurIPS~2025) adapts the Zettelkasten note-taking methodology to agent
memory, treating each stored memory as an evolving \emph{note} in a
dynamic, self-organizing network.  When a new memory is inserted, the
agent autonomously generates a comprehensive note comprising a contextual
description, keywords, and tags; the system then analyzes the existing
memory network to identify relevant historical connections and
establishes typed links between the new note and related prior memories.
Critically, insertion can trigger \emph{retroactive evolution}: existing
memory attributes and context representations are updated when new
information changes their relevance or interpretation, allowing the
network to refine its structure continuously rather than treating
memories as immutable once stored.  Retrieval exploits the link
structure, following chains of association from seed memories to
contextually related nodes rather than performing independent per-memory
scoring.  A-MEM was evaluated across six foundation models and
demonstrated superior performance over state-of-the-art baselines;
it has been adopted as a standard comparative baseline in subsequent
systems (MAGMA reports an A-MEM LLM-as-Judge score of 0.48 on
LoCoMo vs.\ its own 0.70).  The architecture's key contribution is
combining structured organization with agent-driven adaptability:
prior systems either imposed fixed schemas or relied on pure similarity
search; A-MEM's evolving link network provides dynamic structure without
manual curation.  The system does not implement forgetting curves or
preference learning.
